\documentclass[11pt,a4paper]{report}

\usepackage[utf8]{inputenc}
\usepackage[T1]{fontenc}
\usepackage{geometry}
\geometry{margin=1in}
\usepackage{hyperref}
\usepackage{xcolor}
\usepackage{listings}
\usepackage{graphicx}
\usepackage{amsmath}
\usepackage{titlesec}
\usepackage{booktabs}
\usepackage{enumitem}

% Define colors for code listings
\definecolor{codegreen}{rgb}{0,0.6,0}
\definecolor{codegray}{rgb}{0.5,0.5,0.5}
\definecolor{codepurple}{rgb}{0.58,0,0.82}
\definecolor{backcolour}{rgb}{0.95,0.95,0.92}

% Setup listings
\lstset{
    backgroundcolor=\color{backcolour},   
    commentstyle=\color{codegreen},
    keywordstyle=\color{magenta},
    numberstyle=\tiny\color{codegray},
    stringstyle=\color{codepurple},
    basicstyle=\ttfamily\footnotesize,
    breakatwhitespace=false,         
    breaklines=true,                 
    captionpos=b,                    
    keepspaces=true,                 
    numbers=left,                    
    numbersep=5pt,                  
    showspaces=false,                
    showstringspaces=false,
    showtabs=false,                  
    tabsize=2,
    frame=single
}

\title{Building a Production AI Agent from Scratch \\ \large Local File Reader \& Summarizer MVP}
\author{Software Engineering Team}
\date{May 2026}

\begin{document}

\maketitle

\tableofcontents

\chapter{Introduction \& Fundamentals}

\section{LLM Fundamentals}
Large Language Models (LLMs) like Claude are fundamentally \textbf{stateless}. They do not retain memory between independent API calls. To maintain a conversation, the entire history must be sent back to the model with every new request.

In the Anthropic SDK, this is managed via the \texttt{messages} list, which alternates between roles:
\begin{itemize}
    \item \texttt{user}: The human participant.
    \item \texttt{assistant}: The Claude model.
\end{itemize}

\section{The ReAct Loop}
ReAct (Reasoning and Acting) is a paradigm that allows LLMs to interact with external tools. The process follows a loop:
\begin{enumerate}
    \item \textbf{Reasoning}: The model analyzes the user request and decides whether it needs external information.
    \item \textbf{Acting}: If the model decides to use a tool, it returns a \texttt{stop\_reason == "tool\_use"}.
    \item \textbf{Execution}: The application executes the requested Python function.
    \item \textbf{Observation}: The application feeds the result back to the model as a \texttt{tool\_result} within a \texttt{user} turn.
    \item \textbf{Iteration}: The loop repeats until the model provides a final answer (\texttt{stop\_reason == "end\_turn"}).
\end{enumerate}

\section{Function Calling}
Function calling (or Tool Use) involves providing Claude with a JSON Schema that describes a Python function's name, purpose, and parameters. Claude uses this schema to format a request for your code to execute.

\chapter{Project Architecture}

\section{Clean Architecture Layers}
The project follows a 4-layer Clean Architecture to ensure separation of concerns:
\begin{enumerate}
    \item \textbf{Domain Layer}: Core business entities (dataclasses and exceptions) with zero external dependencies.
    \item \textbf{Infrastructure Layer}: Real-world interactions (file system access, session stores).
    \item \textbf{Application Layer}: Orchestrates the business logic and the agentic loop.
    \item \textbf{Presentation Layer}: FastAPI routers and Pydantic schemas for HTTP communication.
\end{enumerate}

\section{Dependency Management}
The project uses \texttt{uv}, a fast Python package manager. The \texttt{pyproject.toml} defines dependencies including \texttt{fastapi}, \texttt{anthropic}, and \texttt{pydantic-settings}.

\chapter{Core Implementation}

\section{The Domain Model}
The \texttt{AgentSession} dataclass represents a conversation thread.
\begin{lstlisting}[language=Python, caption=app/domain/models.py]
@dataclass
class AgentSession:
    session_id: str
    messages: list[dict[str, Any]] = field(default_factory=list)

    def add_user_message(self, content: str) -> None:
        self.messages.append({"role": "user", "content": content})

    def add_assistant_message(self, content: Any) -> None:
        self.messages.append({"role": "assistant", "content": content})

    def add_tool_results(self, tool_results: list[dict[str, Any]]) -> None:
        self.messages.append({"role": "user", "content": tool_results})
\end{lstlisting}

\section{The Agentic Loop}
The \texttt{AgentService} manages the iterative tool-use cycle.
\begin{lstlisting}[language=Python, caption=app/services/agent.py (Simplified)]
async def run(self, session: AgentSession) -> str:
    for iteration in range(MAX_LOOP_ITERATIONS):
        response = await client.messages.create(
            model=settings.claude_model,
            tools=TOOL_DEFINITIONS,
            messages=session.messages,
        )

        if response.stop_reason == "end_turn":
            session.add_assistant_message(serialize(response.content))
            return extract_text(response)

        if response.stop_reason == "tool_use":
            session.add_assistant_message(serialize(response.content))
            tool_results = []
            for block in response.content:
                if block.type == "tool_use":
                    result = execute_tool(block.name, block.input)
                    tool_results.append({
                        "type": "tool_result",
                        "tool_use_id": block.id,
                        "content": result,
                    })
            session.add_tool_results(tool_results)
\end{lstlisting}

\chapter{Security \& Security Patterns}

\section{Path Traversal Prevention}
To prevent unauthorized file access, the tool layer uses \texttt{Path.resolve()} followed by a parent check.
\begin{lstlisting}[language=Python]
requested = (DOCS_DIR / file_name).resolve()
if DOCS_DIR not in requested.parents and requested != DOCS_DIR:
    return "Error: Access denied"
\end{lstlisting}

\section{Secret Management}
API keys are protected using Pydantic's \texttt{SecretStr}, ensuring they are never logged or exposed in JSON outputs.

\chapter{Conclusion}
This architecture provides a robust, production-ready foundation for AI agents. By decoupling the LLM loop from the transport layer and enforcing strict security patterns, we build agents that are both capable and safe.

\end{document}
